% SPDX-License-Identifier: CC-BY-NC-ND-4.0
% !TEX root = main.tex

\section{要素の圏}

\begin{definition}
$\mathscr{C}$を局所小圏とし$F \colon \mathscr{C} \to \mathbf{Set}$を共変関手とする.\ $A \in \Ob(\mathscr{C}),\ x \in F(A)$としたとき対象を$(A, x)$,\ 射$(A, x) \to (B, y)$を$\mathscr{C}$の射$f \colon A \to B$で$F(f)(x) = y$を満たすもの,\ 射の合成を$\mathscr{C}$の射の合成で定めることで圏$\displaystyle \int_{\mathscr{C}} F$が定まる.\ これを$F$の\textbf{要素の圏}(category of elements)という.\ 反変関手$F \colon \mathscr{C} \to \mathbf{Set}$に対しては$A \in \Ob(\mathscr{C}),\ x \in F(A)$としたとき対象を$(A, x)$,\ 射$(A, x) \to (B, y)$を$\mathscr{C}$の射$f \colon A \to B$で$F(f)(y) = x$を満たすものとして要素の圏を定める.
\end{definition}

\begin{definition}
$\mathscr{C}$を圏とする.\ $I \in \Ob(\mathscr{C})$は任意の$X \in \Ob(\mathscr{C})$に対してただ1つ$f \in \Mor_\mathscr{C}(I,X)$が存在するとき\textbf{始対象}(initial object)であるという.
\end{definition}

\begin{definition}
$\mathscr{C}$を圏とする.\ $T \in \Ob(\mathscr{C})$は任意の$X \in \Ob(\mathscr{C})$に対してただ1つ$f \in \Mor_\mathscr{C}(X, T)$が存在するとき\textbf{終対象}(terminal object)であるという.
\end{definition}

\begin{definition}
$\mathscr{C}$を圏とする.\ $Z \in \Ob(\mathscr{C})$は始対象かつ終対象であるとき\textbf{零対象}(zero object)であるという.
\end{definition}

\begin{proposition}
$\mathscr{C}$を局所小圏,\ $F \colon \mathscr{C} \to \mathbf{Set}$を共変関手,\ $A \in \Ob(\mathscr{C}), u \in F(A)$とする.\ このとき次が成り立つ.
\begin{center}
$(A, u)$は$F$の普遍対$\iff$$(A,u)$は$\displaystyle \int_\mathscr{C} F$の始対象
\end{center}
\end{proposition}

\begin{proof}
任意の$(B,x) \in \Ob(\int_\mathscr{C} F)$に対して$\Mor_{\int_\mathscr{C} F}((A, u), (B, x)) = \{f \in \Mor_\mathscr{C}(A, B) \mid F(f)(u) = x\}$であり$(A, u)$が始対象であることと任意の$(B,x) \in \Ob(\int_\mathscr{C} F)$に対して$\Mor_{\int_\mathscr{C} F}((A, u), (B, x))$がシングルトンであることは同値である.\ そして任意の$(B,x) \in \Ob(\int_\mathscr{C} F)$に対して$\Mor_{\int_\mathscr{C} F}((A, u), (B, x))$がシングルトンであることは\cref{def: IsUniversalElement}の(\ast)と同値である.
\end{proof}